\documentclass[10pt]{article}
\usepackage[margin=0.6in]{geometry}
\usepackage{enumitem}
\usepackage{booktabs}
\usepackage{graphicx}
\usepackage{xcolor}
\usepackage{hyperref}
\usepackage{fancyhdr}

\pagestyle{fancy}
\fancyhf{}
\renewcommand{\headrulewidth}{0pt}
\fancyfoot[C]{\footnotesize Akum Gill $\cdot$ \today}

\setlength{\parindent}{0pt}
\setlength{\parskip}{4pt}

\begin{document}

\begin{center}
{\Large \textbf{MAGNETS Spectroscopic Orchestration Layer}}\\[4pt]
{\normalsize Infrastructure for Queue-Based LLAMAS Observing}
\end{center}

\vspace{-8pt}
\hrule
\vspace{8pt}

\textbf{Context.} The MAGNETS collaboration pools Magellan time across partner institutions for Rubin transient follow-up. Per the Stubbs 2026B proposal: ``We will apportion queue observing time and target prioritization in proportion to time awarded by respective TACs.'' This requires automated infrastructure to manage the queue schedule.

\textbf{Problem.} Current workflow (via Ashley's group) uses manual notebook-based scheduling. This doesn't scale with multiple requesters, lacks time accounting, and requires format translation between request intake and scheduling scripts.

\textbf{Proposed Solution.} An orchestration layer connecting RubinAlerts candidate discovery to LLAMAS observing plan generation:

\vspace{-4pt}
\begin{center}
\small
\begin{tabular}{@{}p{2.8cm}p{10cm}@{}}
\toprule
\textbf{Component} & \textbf{Function} \\
\midrule
Ingester & Pull targets from RubinAlerts \texttt{candidates.csv} and manual requests \\
Normalizer & Convert to common schema; estimate exposure times by redshift \\
Prioritizer & Score by merit $\times$ observability $\times$ moon phase compatibility \\
Time Accountant & Track 30-hour allocation across moon phases (5D + 20G + 5B) \\
LLAMAS Planner & Greedy scheduling with WD standard interleaving \\
\bottomrule
\end{tabular}
\end{center}

\textbf{LLAMAS-Specific Scheduling.} From proposal Table 1, exposure times scale with redshift and moon phase:

\vspace{-4pt}
\begin{center}
\small
\begin{tabular}{@{}ccccc@{}}
\toprule
\textbf{Redshift} & \textbf{Moon} & \textbf{Peak $r$} & \textbf{Exposure} & \textbf{Binning} \\
\midrule
0.10--0.20 & Bright & 19.3 & 35 min & 2$\times$ \\
0.20--0.30 & Grey & 20.6 & 95 min & 2$\times$ \\
0.30--0.35 & Grey/Dark & 21.3 & 45--160 min & 4$\times$ \\
0.35--0.40 & Dark & 21.9 & 180 min & 4$\times$ \\
\bottomrule
\end{tabular}
\end{center}

\vspace{-2pt}
IFU advantages: $\sim$1 min overhead (vs.\ $\sim$10 min for slit), no slit losses, host galaxy ``for free.''

\textbf{Integration with RubinAlerts.} The existing pipeline produces nightly \texttt{candidates.csv} with merit-scored SNe~Ia. The orchestration layer will:
\begin{itemize}[nosep, leftmargin=1.5em]
    \item Filter to $z = 0.1$--$0.4$, $r = 18$--$21.5$, elliptical hosts preferred
    \item Map merit scores to LLAMAS priorities
    \item Generate queue-ready observing plans with WD standards (Boyd et al.\ 2026)
\end{itemize}

\textbf{Implementation Timeline.}
Phase 1 (Wk 1--2): Normalizer, LLAMAS planner, CLI.
Phase 2 (Wk 3): Time accounting, prioritizer.
Phase 3 (Wk 4): RubinAlerts integration.
Phase 4 (Wk 5+): WD standards, reporting.

\textbf{Open Questions.}
(1) LLAMAS-specific constraints to confirm with Rob Simcoe;
(2) Time charging: on schedule or observation?
(3) Multi-night lookahead vs.\ nightly optimization?

\textbf{Resources.} Full design doc: \texttt{docs/design/spectroscopic-orchestration.md}

\end{document}
